\chapter{Przedstawienie problemu harmonogramowania zabiegów rehabilitacyjnych}
Ten rozdział ma na celu przybliżenie specyfiki problemu harmonogramowania zabiegów
rehabilitacyjnych, jaki występuje w placówkach współpracujących z firmą Medinet.

\section{Podstawowy problem harmonogramowania}
Problem harmonogramowania zabiegów rehabilitacyjnych jest wariantem ogólnego problemu
harmonogramowania (z ang. \emph{timetable scheduling}). Problem ten  stara się
przypisać pewnym wydarzeniom - kursom, zabiegom, lotom itd. - pewien moment
wykonania (termin) oraz pewne niezbędne zasoby - sale, nauczycieli,
rehabilitantów, sprzęt itd. - tak, aby spełnić pewne ograniczenia oraz
zmaksymalizować poziom zaspokojenia określonych preferencji. Ograniczenia muszą
zostać spełnione aby rozwiązanie było dopuszczalne, natomiast preferencje muszą
być zaspokojone w największym możliwym stopniu.

Ogólny problem harmonogramowania można zapisać \cite{habib2013} jako:
zbiór terminów \emph{T}, zbiór zasobów \emph{R}, zbiór spotkań \emph{M} oraz
zbiór ograniczeń \emph{C}. Rozwiązaniem podstawowego problemu harmonogramowania,
jest takie przypisanie terminów oraz niezbędnych zasobów do spotkań, aby jak
najlepiej spełnić konkretne ograniczenia. Wariant podstawowy nie definiuje
osobno preferencji i ograniczeń twardych, tylko chęć spełnienia jak największej ilości ograniczeń
jako takich. To czy dane ograniczenie jest wymagane czy też nie, to szczegół
implementacyjny uzależniony od specyfiki konkretnej wersji problemu.

Jak się okaże w kolejnym rozdziale, problem harmonogramowania zabiegów
rehabilitacyjnych jest de facto nieco bardziej skomplikowaną wersją ogólnego
problemu harmonogramowania.
\pagebreak
\section{Definicja problemu rehabilitacyjnego}
Placówki świadczące usługi rehabilitacyjne operują na tzw. zleceniach. Zlecenie
\sRef\ jest to lista pomniejszych zleceń \sSubRef\ na pewien konkretny zabieg rehabilitacyjny.
Każdy zabieg trwa określoną ilość czasu oraz jest odbywany na konkretnej sali
zabiegowej. Zadaniem pracownika jest takie przypisanie terminów oraz sal do konkretnych
zabiegów, aby jak najlepiej spełnić ograniczenia nadane przez placówkę.

\begin{figure}[h]
	\centering
	\begin{tabular}{ l l l l }
	Id zlecenia & Id zabiegu & Nazwa zabiegu & Czas trwania zabiegu \\
	\hline
	14123 & 0 & Laseroterapia & 15m \\
	14124 & 0 & Laseroterapia & 15m \\
	14125 & 2 & Krioterapia & 7m \\
	\end{tabular}
	\caption{Przykładowe zlecenie \sRef}
	\label{example-referral}
\end{figure}

Problem można zdefiniować jako: \emph{$SP=\{\sRef, start, end\}$}. Gdzie \sRef\
jest to lista zabiegów, \emph{start} jest datą początkową a \emph{end} końcową w
obrębie których mają się znaleźć wyznaczone terminy dla zabiegów z listy \sRef.
Rozwiązaniem problemu \emph{SP} jest
\emph{$\sSol=\{\sTermSet_1,\sTermSet_2...\sTermSet_n\}$}, gdzie \sSol\ jest
zbiorem pomniejszych zbiorów \sTermSet, które zawierają w sobie konkretne terminy
\sTerm. Pozycja zbiorów \sTermSet\ określa to któremu zabiegowi odpowiadają dane
terminy. Trzymając się przykładu z rysunku \ref{example-referral}:

\begin{figure}[h]
	\centering
	\begin{tabular}{ l l l l }
	Id zlecenia & Id terminu & Początek terminu & Czas trwania terminu \\
	\hline
	14123 & 645 & 08:00 01.01.2021 & 10m \\
	14123 & 646 & 08:10 01.01.2021 & 10m \\
	14124 & 816 & 10:00 02.01.2021 & 20m \\
	14125 & 910 & 08:30 01.01.2021 & 15m \\
	\end{tabular}
	\caption{Przykładowe rozwiązanie \emph{S}}
	\label{example-solution}
\end{figure}

Zauważmy, że pierwsze dwa terminy są przypisane do tego samego zlecenia,
ponieważ czas trwania terminu jest mniejszy niż czas trwania zabiegu.

To jakie terminy \sTerm\ mogą zostać przypisane danemu zleceniu \sSubRef,\
zależy od typu zabiegu. Każdy typ zabiegu np. laseroterapia, masaż etc. jest
połączony z odpowiednimi \emph{listami terminów}. Listy terminów określają
konkretne terminy na dany dzień dla danej sali zabiegowej \sPlace, a sama sala
ogranicza nam zabiegi, które mogą się w niej wykonywać. Można śmiało myśląc o
listach terminów, myśleć o konkretnych dniach miesiąca. Przykładowo lista na
dzień 01.01.2021 dla sali "Basen" może zawierać w sobie 30 terminów od godziny 08:00
do 18:00 w interwałach 20 minutowych. Czas trwania terminu dla tej listy wynosi
więc 20 minut. Terminy w zestawie mogą mieć różną długość: 5m, 10m, 15m, 20m, 30m
lub 60m. % TODO: Sprawdź konkretne czasy

\pagebreak
\begin{figure}[h]
	\centering
	\includegraphics[width=\textwidth]{gfx/medinet-treatments-stats.pdf}
	\caption{Ilość zabiegów w zleceniu}
	\label{medinet-treatments-stats}
\end{figure}
Z bazy danych firmy Medinet możemy uzyskać następujące dane statystyczne co do
ilości zabiegów w zleceniach. Jak widać ilość zabiegów wacha się między 10 a 70.
Najczęstsza ilość zabiegów w zleceniu to 30, później 40 itd. Duża powtarzalność
ilości zabiegów jest spowodowana szablonami, z których korzystają lekarze.
Szablony nie są ściśle określone przez prawo ani nawet przez jednostkę, lecz są
to pomocnicze elementy, które zależą od konkretnego lekarza i jego praktyki.
\pagebreak
\section{Ograniczenia}
Wiemy już jak wygląda problem oraz jego rozwiązanie. Specyfika problemu
harmonogramowania zabiegów rehabilitacyjnych wymusza wzięcie pod uwagę
dodatkowych ograniczeń, które musi spełniać rozwiązanie \sSol.
\begin{enumerate}
	\item{$C_1$ - Dwa zabiegi rehabilitacyjne tego samego typu nie mogą być
		zaplanowane na ten sam dzień.}
	\item{$C_2$ - Dwa zabiegi rehabilitacyjne nie mogą być zaplanowane w na ten sam czas.}
	\item{$C_3$ - Zabiegi określone w słowniku \sGap\ muszą zachować określony tam odstęp.}
	\item{$C_4$ - Zabiegi nie mogą zostać zapisane w terminach, które są zapełnione.}
\end{enumerate}
Słownik \sGap\ jest to struktura przechowująca informację o minimalnych
przerwach między danymi typami zabiegów.
Jak widać na poniższych rysunkach, zabiegi mogą mieć zdefiniowaną różną przerwę
w zależności od tego, który zabieg występuje jako pierwszy w danym dniu. Kolorem
czerwonym została zaznaczona wymagana przerwa, natomiast zabieg 0 i 1
reprezentowane są przez odpowiednio oznaczone prostokąty.
\\
\begin{figure}[h]
	\centering
	\begin{tabular}{ l l l l l }
		Id zabiegu & Id następnego zabiegu & Wymagana przerwa \\
	\hline
		1 & 0 & 60m \\
		0 & 1 & 20m \\
	\end{tabular}
	\caption{Przykładowa struktura \sGap\ definiująca wymagane przerwy}
	\label{example-gaps}
\end{figure}
\\
\begin{figure}[h]
	\centering
	\includegraphics[width=\textwidth]{gfx/gap_1.pdf}
	\caption{Wizualizacja wymaganej przerwy między zabiegami (A)}
	\label{example-gaps-1}
\end{figure}
\\
\begin{figure}[h]
	\centering
	\includegraphics[width=\textwidth]{gfx/gap_2.pdf}
	\caption{Wizualizacja wymaganej przerwy między zabiegami (B)}
	\label{example-gaps-2}
\end{figure}
\pagebreak
\section{Preferencje}
Preferencje są to cechy, które powinno spełniać dobre rozwiązanie.
Aczkolwiek ich spełnienie nie jest wymagane aby rozwiązanie było akceptowalne.
\begin{enumerate}
	\item{$P_1$ - Im wcześniej pacjent opuści placówkę tym lepiej. }
	\item{$P_2$ - Im mniej przerw między zabiegami w ciągu dnia tym lepiej.}
	\item{$P_3$ - Im mniej wykorzystanych dni tym lepiej.}
	\item{$P_4$ - Im wcześniej zaczyna się dany zabieg w danym dniu tym lepiej.}
	\item{$P_5$ - Im większa ilość zabiegów w danym dniu tym lepiej.}
\end{enumerate}

Jak widać preferencję dotyczą głównie chęci minimalizacji czasu spędzonego na
obsłudze danego zlecenia. Im mniej zasobów jednostki jest wykorzystanych tym
lepiej.

Drugorzędnym celem jest maksymalizacja komfortu pacjenta. Im mniej czasu pacjent
musi spędzić na jednostce tym lepiej dla niego jak i dla samej jednostki.
Preferencje czasowe pacjenta nie są uwzględnione w samym wyznaczaniu terminów,
ponieważ priorytetem jest jak najefektywniejsza obsługa zleceń. Z praktyki
jednostek wynika, że preferencje pacjentów są bardzo nieprecyzyjne, mgliste i
nie dają się w konkretny sposób przełożyć na proces wyznaczania.
